\section{Finite-adelically distorted dynamical systems}
\label{sec:dynamics}

A self-map \(f:X\to X\) is \emph{confined} if
\(f_n=\#\Fix(f^n)\) is finite for every \(n\geq1\).
Its fixed-count series and dynamical zeta function are
\[
Z_f(z)=\sum_{n\geq1}f_nz^n,\qquad
\zeta_f(z)=\exp\left(\sum_{n\geq1}\frac{f_n}{n}z^n\right).
\]
A sequence of period \(m\) is a \emph{gcd sequence} if its value at \(n\)
depends only on \(\gcd(n,m)\). We use the following finite-adelically
distorted, or FAD, presentation from
\citep[Definitions 7.1.1--7.1.2]{bch2024}:
\begin{equation}\label{eq:fad}
f_n=|\det(A^n-I)|c^nr_n
\prod_{p\in S}|n|_p^{s_{p,n}}
p^{-t_{p,n}|n|_p^{-1}}.
\end{equation}
Here \(A\) is an integral square matrix, \(c>0\), \(S\) is finite,
\(r_n\) is a positive real gcd sequence, and \(s_{p,n},t_{p,n}\)
are nonnegative real gcd sequences having periods coprime to their
own prime \(p\). A prime is \emph{genuinely distorted} if its exponent
pair is active at some residue. An exponent \(t_{p,n}\ne0\) is the
wild part of the distortion.

We use counting entropy, namely
\(\limsup_{n\to\infty}n^{-1}\log f_n=\log\Lambda\) for a nonzero
count sequence. The standard FAD entropy formula identifies
\(\Lambda=c\prod_{|\lambda|>1}|\lambda|\), with the eigenvalues of
\(A\) counted with multiplicity
\citep[Section 10.3]{bch2024}.
Our argument below also identifies the radius \(\Lambda^{-1}\)
directly whenever distortion is present.

\begin{corollary}\label{cor:fad}
Let a nonzero confined FAD system have positive counting entropy
\(\log\Lambda>0\). If at least one prime in \eqref{eq:fad} is genuinely
distorted, then both \(Z_f\) and \(\zeta_f\) have
\(|z|=\Lambda^{-1}\) as a natural boundary, with no meromorphic
continuation through any point of that circle.
No unique-dominant-root hypothesis or integrality assumption on
the wild exponents is required.
\end{corollary}

\begin{proof}
First, \(A\) has no root-of-unity eigenvalue. If such an eigenvalue
had order \(h\), then \(f_n=0\) at every multiple of \(h\).
A positive \(f_m\) would supply a point fixed also by \(f^{hm}\),
a contradiction. The assumed nonzero count sequence therefore
excludes this possibility.

The determinant factor has an expansion
\begin{equation}\label{eq:leading}
|\det(A^n-I)|c^n=\Lambda^n u_n+O((\Lambda\tau)^n),
\qquad 0\leq\tau<1,
\end{equation}
where \(u_n\) is a finite sum of unit phases and is strictly positive
for every positive integer. This is the classical dominant expansion
in \citep[Section 10.3, especially Lemma 10.3.10]{bch2024}.
Its relevant features can also be seen directly. After the moduli of
the eigenvalues outside the unit circle have been extracted, the
off-circle factors tend to one exponentially. Unit-circle eigenvalues
occur in nonreal conjugate pairs, since the real possibilities \(1\)
and \(-1\) have been excluded. Their product contributes
\[
u_n=\prod_{\{\eta,\bar\eta\}}
(2-\eta^n-\eta^{-n})>0.
\]
Expanding this finite product gives the required finite phase sum;
the empty product is one.

Let \(p\) be a genuinely distorted prime, and choose an active
exponent residue \(n_0\bmod m_p\), where \(m_p\) is a common period of
its two exponent sequences and \(p\nmid m_p\). Choose \(W\) for the
dominant coefficient series as in Section~\ref{sec:criterion}.
Since \(m_p\mid W\), the Chinese remainder theorem supplies
\(a\bmod W\) satisfying
\[
a\equiv0\pmod{p^{v_p(W)}},
\qquad a\equiv n_0\pmod{m_p}.
\]
The exponent pair is active at \(a\). If \(b_C(a)=0\) for every phase
class, positivity of \(r_a\) would make every grouped periodic phase
sum vanish there. Consequently \(u_n\) would vanish for all positive
\(n\equiv a\pmod W\), contradicting \(u_n>0\).
Thus \(\AF\) holds.

Apply Theorem~\ref{thm:main} to the normalized leading part of
\eqref{eq:leading}, multiplied by the periodic and local factors.
It has dense nonzero radial masses on the unit circle.
Those factors are bounded: the local kernels are at most one and
\(r_n\) has a finite period. Hence the normalized error in
\eqref{eq:leading} generates a function holomorphic in a disc of
radius greater than one. Adding it does not change any nonzero
radial mass or remove the resulting dense singularities. Rescaling
proves the assertion for \(Z_f\).

Within \(|z|<\Lambda^{-1}\), the defining exponential for \(\zeta_f\)
converges and is nonvanishing, and
\begin{equation}\label{eq:log-derivative}
z\frac{\zeta_f'(z)}{\zeta_f(z)}=Z_f(z).
\end{equation}
If \(\zeta_f\) had a meromorphic continuation through a boundary point,
it would not be identically zero in the continued neighborhood, because
it agrees with the nonvanishing exponential on the interior overlap.
Its zeros and poles would be discrete. A smaller nonempty boundary
arc would therefore avoid them, and \eqref{eq:log-derivative} would
continue \(Z_f\) holomorphically through that arc. This contradicts
the dense singularities already obtained.
\end{proof}

The own-prime coprimality of the exponent periods has one precise
function in this argument: it permits an active exponent value and an
unbounded valuation fibre to occur together. Positivity of the leading
factor then prevents all phase classes from being masked.
Neither hypothesis is needed in the abstract criterion, which retains
the rational masked cases.

\subsection{The already-known no-wild deduction}

We record why removing unique dominance only in the no-wild FAD case
is not an additional contribution. Suppose every \(t_{p,n}=0\), and
write \(f_n=A_ng(n)\), where
\[
A_n=|\det(A^n-I)|c^n,\qquad
g(n)=r_n\prod_{p\in S}p^{-s_{p,n}v_p(n)}.
\]
For an actual FAD system, \citep[Proposition 10.2.1]{bch2024} gives
\(c\in\Z_{>0}\). Since the wild exponents are zero and hence integral,
\citep[Proposition 10.2.2]{bch2024} gives rationality of \(r_n\) and
\(p^{s_{p,n}}\). Therefore \(g(n)\in\Q\), and its logarithmic Weil height
satisfies \(h(g(n))=O(\log n)\).

If \(S\) is empty, \(g\) is periodic and already an almost
quasi-polynomial. Otherwise, let \(W\) be a common period and put
\(d_k=W\prod_{p\in S}p^k\).
On each residue class modulo \(d_k\) which avoids zero modulo
\(p^k\) for every \(p\in S\), all relevant valuations and periodic
data are fixed. Thus \(g\) is constant there. The proportion of
excluded classes is at most \(\sum_{p\in S}p^{-k}\), tending to zero.
These rationality, height, and density-one local polynomial conditions
make \(g\) an almost quasi-polynomial in the sense of
\citep[Definition 2.11]{bhn2025}, with constant constituent polynomials.

The generating function of \(A_n\) is rational: expansion of the
determinant gives finitely many exponential terms, and its real
sign is periodic. Its dominant part is strictly positive by
\eqref{eq:leading}, so it vanishes identically on no arithmetic
progression. It is therefore stable in the sense required by
\citep[Theorem 2.14]{bhn2025}. That theorem yields the rational or
natural-boundary alternative for \(\sum f_nz^n\).
If a valuation exponent is active, the FAD rationality criterion
\citep[Proposition 11.3.3]{bch2024} excludes rationality.
The logarithmic derivative then transfers the boundary to
\(\zeta_f\). In particular, the abelian-variety and finite-place
solenoid applications already settled in \citep{bhn2025} are not
being asserted anew here.

The classical root-rational classification for FAD zeta functions
\citep[Theorem 11.3.4]{bch2024} is a separate input and is not reproved
as part of Theorem~\ref{thm:main}. Our additional analytic argument
is needed for wild data beyond the preceding height hypothesis,
as the next example makes explicit.
